\documentclass{article}
\usepackage{booktabs}
\usepackage{amsmath}
\usepackage{microtype}

% Configuramos la hoja en orientación horizontal (landscape)
\usepackage[landscape, left=2cm, right=2cm, top=2cm, bottom=2cm]{geometry}

\begin{document}

\begin{table}[htbp]
\centering
% Reducimos un poco el tamaño de fuente y el espacio entre columnas para máxima legibilidad
\small
\setlength{\tabcolsep}{5pt}
\renewcommand{\arraystretch}{1.15} % Le da un poco de aire vertical a las filas

\caption{Avance 5: Resultados de rendimiento y comparación de métodos de mejora (Enhancement) frente al baseline, usando pesos re-entrenados en SCARED (Dr. Ricardo Espinosa Loera).}
\label{tab:resultados_avance5}

\begin{tabular}{llccccccccc@{\hspace{12pt}}l}
\toprule
\textbf{Modelo} & \textbf{Enhancement} & \textbf{AbsRel} & \textbf{SqRel} & \textbf{RMSE} & \textbf{RMSELog} & $\mathbf{\delta < 1.25}$ & $\mathbf{\delta < 1.25^2}$ & $\mathbf{\delta < 1.25^3}$ & \textbf{IT (ms)} & \textbf{$\mathbf{\Delta\text{AR}\%}$} & \\
\midrule

Monodepth2 & none       & 0.0568 & 0.6821 &  4.771 & 0.0777 & 0.9852 & 0.9915 & 0.9939 &  27.4 & baseline  & \\
           & retinex    & 0.1114 & 1.4798 &  8.236 & 0.1489 & 0.8656 & 0.9742 & 0.9897 & 413.0 & +96.1\%   & \\
           & endolmspec & 0.0662 & 0.8576 &  5.324 & 0.0876 & 0.9787 & 0.9912 & 0.9929 & 150.0 & +16.5\%   & \\
           & iat        & 0.0581 & 0.6269 &  4.796 & 0.0791 & 0.9832 & 0.9913 & 0.9948 & 285.3 &  +2.3\%   & \\
\addlinespace[0.8em]

MonoViT & none       & 0.1328 & 1.9229 &  9.966 & 0.1752 & 0.8476 & 0.9685 & 0.9868 &  51.4 & baseline  & \\
        & retinex    & 0.1339 & 1.9151 &  9.809 & 0.1755 & 0.8397 & 0.9535 & 0.9887 & 437.5 &  +0.8\%   & \\
        & endolmspec & 0.1425 & 2.1251 & 10.541 & 0.1841 & 0.8173 & 0.9666 & 0.9864 & 173.9 &  +7.3\%   & \\
        & iat        & 0.1353 & 1.9456 &  9.930 & 0.1758 & 0.8464 & 0.9643 & 0.9874 & 309.4 &  +1.9\%   & \\
\addlinespace[0.8em]

EndoSfMLearner & none       & 0.1706 & 3.1434 & 12.754 & 0.2074 & 0.7750 & 0.9387 & 0.9884 &   6.4 & baseline  & \\
               & retinex    & 0.1710 & 3.1565 & 12.777 & 0.2077 & 0.7743 & 0.9386 & 0.9884 & 392.2 &  +0.2\%   & \\
               & endolmspec & 0.1707 & 3.1466 & 12.761 & 0.2074 & 0.7747 & 0.9387 & 0.9884 & 129.1 &  +0.1\%   & \\
               & iat        & 0.1707 & 3.1476 & 12.761 & 0.2075 & 0.7747 & 0.9387 & 0.9884 & 264.3 &  +0.1\%   & \\
\addlinespace[0.8em]

AF-SfMLearner & none       & 0.1103 & 1.3733 & 8.970 & 0.1420 & 0.9008 & 0.9870 & 0.9966 &  27.2 & baseline  & \\
              & retinex    & 0.1326 & 1.8277 & 9.682 & 0.1786 & 0.8552 & 0.9600 & 0.9874 & 413.1 & +20.2\%   & \\
              & endolmspec & 0.1055 & 1.2379 & 8.483 & 0.1365 & 0.9179 & 0.9869 & 0.9968 & 150.2 &  -4.4\%   & $\leftarrow$ \\
              & iat        & 0.1079 & 1.2953 & 8.678 & 0.1408 & 0.9026 & 0.9844 & 0.9967 & 285.3 &  -2.2\%   & $\leftarrow$ \\
\addlinespace[0.8em]

MonoIIT & none       & 0.0373 & 0.2642 & 3.487 & 0.0572 & 0.9886 & 0.9958 & 0.9995 &  72.7 & baseline  & \\
        & retinex    & 0.0988 & 1.1629 & 7.732 & 0.1353 & 0.8911 & 0.9822 & 0.9947 & 458.5 & +164.9\%  & \\
        & endolmspec & 0.0424 & 0.3540 & 3.700 & 0.0618 & 0.9817 & 0.9926 & 0.9985 & 195.5 & +13.7\%   & \\
        & iat        & 0.0411 & 0.3067 & 3.637 & 0.0600 & 0.9848 & 0.9941 & 0.9994 & 330.7 & +10.2\%   & \\
\addlinespace[0.8em]

w19-MonoViT & none       & 0.0459 & 0.2755 & 3.913 & 0.0637 & 0.9881 & 0.9986 & 0.9997 &  72.5 & baseline  & \\
            & retinex    & 0.1342 & 2.2202 & 9.920 & 0.1850 & 0.8114 & 0.9305 & 0.9740 & 458.4 & +192.4\%  & \\
            & endolmspec & 0.0597 & 0.3932 & 4.705 & 0.0786 & 0.9827 & 0.9984 & 0.9997 & 195.4 & +30.1\%   & \\
            & iat        & 0.0524 & 0.3212 & 4.108 & 0.0689 & 0.9880 & 0.9978 & 0.9996 & 330.6 & +14.2\%   & \\
\bottomrule
\end{tabular}

\vspace{0.6em}
\begin{minipage}{\linewidth}
    \footnotesize
    \textbf{Nota:} Pesos re-entrenados en SCARED. $\delta < 1.25^k$ representa la fracción de píxeles con $\max(\text{pred}/\text{gt}, \text{gt}/\text{pred}) < 1.25^k$ (valores más altos indican un mejor rendimiento).
    Un porcentaje negativo en $\Delta\text{AR}\%$ indica una mejora (\textit{improvement}) respecto a la configuración sin procesar (\textit{none}). Los casos que logran esta mejora se señalan explícitamente con el símbolo $\leftarrow$.
\end{minipage}

\end{table}

\end{document}
